%!TEX root = TQFTbook.tex
%%%%%%%%%%%%%%
\chapter{Higher categories}\label{c:higher-cat}

In the previous chapters, we have defined extended 2d TQFTs using
the language of weak 2-categories and symmetric monoidal 2-categories. In a
similar way, it is clear that to define extended TQFTs in higher
dimensions, we will need higher categories: if we have a good definition of
an $n$-category, and we can construct the category of cobordisms with
corners $\Bord_n$, we can define an extended $n$-dimensional TQFT as a
symmetric monoidal functor $\Bord_n\to \calC$, for some symmetric monoidal
$n$-category $\calC$.

In this chapter, we discuss various approaches to defining higher categories. We
will restrict ourselves to a very brief overview, referring the reader to
original papers for details.

%%%%%%%%%%%%%%%%
\section{Fundamental groupoid}\label{s:fundamental-groupoid}

Reversing the traditional order in mathematics, we begin our discussion of
higher categories not with a definition but with an example: whatever definition
we give later, it should cover that example. This example is that of a
fundamental groupoid, described informally below.


Recall that for a topological space $X$, we denote by $\Pi(X)$ the fundamental
groupoid of $X$; it is a category whose objects are points of $X$ and morphisms
are homotopy classes of paths in $X$.

We can extend it and define the fundamental infinity groupoid $\Pi_{\infty}(X)$
as follows:
\begin{itemize}
  \item Objects of $\Pi_{\infty}(X)$ are points of $X$.
  \item 1-morphisms $x_0\to x_1$ are paths in $X$, i.e., continuous maps
    $I\to X$ such that $\ga(0)=x_0$, $\ga(1)=x_1$ (here and below, $I=[0,1]$).
  \item 2-morphisms $\ga_0\to \ga_1$, where $\ga_i$ are paths from $x_0$ to
  $x_1$, are path homotopies, i.e., continuous maps $h\colon I^2\to X$ such that
  $h(0,t)=\ga_0(t)$, $h(1, t)=\ga_1(t)$, and for all $s$, $h(s,0)=x_0$,
  $h(s,1)=x_1$.
  \item 3-morphisms are homotopies between homotopies, i.e., continuous maps
  $I^3\to X$ with suitable restrictions to the boundary.
  \item \dots
\end{itemize}
This defines an ``infinity groupoid'', which has morphisms of all orders. We can
also define the fundamental $n$-groupoid $\Pi_{\leq n}(X)=\tau_{\leq
n}(\Pi_\infty(X))$, which only has morphisms up to order $n$: namely,
$n$-morphisms in $\Pi_{\leq n}(X)$ are homotopy classes of
$n$-morphisms in $\Pi_{\infty}(X)$ (two $n$-morphisms are homotopic if there
exists an $(n+1)$-morphism between them in $\Pi_\infty(X)$). For $n=1$, this
coincides with the previous definition of the fundamental groupoid:
$\Pi_{\le 1}(X)=\Pi(X)$.

Note that we have not yet discussed composition in $\Pi_\infty(X)$. It is easy
to define composition in the same way as it is done in the usual fundamental
groupoid $\Pi_{\le 1}(X)$; however, so defined composition is not associative in
$\Pi_\infty$. For example, for $1$-morphisms $\ga_1, \ga_2, \ga_3$, the
compositions $\ga_3 (\ga_2 \ga_1)$ and $(\ga_3 \ga_2)\ga_1$ are not equal as
maps $I\to X$: in the first case, $\ga_1$ is traveled over $t\in [0, 1/4]$,
$\ga_2$ over $t\in [1/4, 1/2]$, and $\ga_3$ over $[1/2,1]$. For the second
composition, the intervals are $[0, 1/2]$, $[1/2, 3/4]$, and $[3/4,1]$
respectively, so these compositions are not equal; however, they are homotopic
to each other. Thus, any definition of an $n$-category should allow for
compositions ``associative up to homotopy''.

Fundamental $n$-groupoids $\Pi_{\leq n}(X)$ were suggested by Grothendieck in
\ocite{grothendieck}. He also proposed the following hypothesis.
\begin{conjecture}[Homotopy hypothesis]
  Homotopy $n$-types up to weak equivalence
  are classified by $\Pi_{\leq n}(X)$ up to (suitable) equivalence.
\end{conjecture}
(Recall that a topological space $X$ is called an \emph{$n$-type} if its
homotopy groups $\pi_k(X)$ are trivial for all $k>n$, and a continuous function
$f\colon X\to Y$ is a weak equivalence if for every $k\ge 0$, $f$ induces an
isomorphism of homotopy groups $\pi_k(X)\cong \pi_k(Y)$.)

Of course, in order to discuss this hypothesis, one needs first to give a
definition of $n$-categories and equivalences between them; thus, we should
consider this hypothesis as a natural requirement which should be satisfied by
any definition of an $n$-category.

One can also state an infinity version of the homotopy hypothesis; namely, it
says that topological spaces up to weak equivalence are classified by their
fundamental infinity groupoids up to equivalence.

%%%%%%%%%%%%%%%%
\section{Strict and weak $n$-categories}\label{s:strict-n-cat}
After giving the motivating example of $n$-groupoids, we can try to
define the notion of an $n$-category.

The simplest approach is to consider strict $n$-categories. These categories are
not very hard to define. Namely, a strict $n$-category must have objects, and
for any objects $X,Y$, we must have an $(n-1)$-category $\Hom (X,Y)$, together
with a composition functor of $(n-1)$-categories $\Hom(X,Y)\times \Hom(Y,Z)\to
\Hom(X,Z)$, which must be associative. In a similar way one defines identity
morphisms. (Category theorists would say that an
$n$-category is a category enriched over $(n-1)$-categories.)

However, this definition is not of much use. First, virtually all known examples
are non-strict --- compositions are not strictly associative but only
``associative up to higher order isomorphisms'', as we have already seen in the
example of fundamental groupoids. Moreover, for $n\ge 3$, it turns out that this
is more than just a technical problem: not every weak 3-category is equivalent
to a strict 3-category (unlike the case of 2-categories). For example, it can be
shown that $\Pi_{\leq 3}(S^2)$, the truncated fundamental
groupoid of $S^2$, is not equivalent to a strict 3-groupoid (see
\ocite{simpson}*{Theorem 2.7.2}).



Thus, we have no choice but to forget strict categories and start developing the
theory of weak $n$-categories, where compositions are ``associative up to
higher order isomorphisms''.


Unfortunately, giving an algebraic definition of a weak $n$-category is very
hard to do. A definition of a weak 3-category exists (see
\ocite{gordon-power-street}), but it is so long that most people give up before
reaching the end of it; for $n\ge 4$, the algebraic definitions are even more
cumbersome. Some summary of algebraic approaches is collected in
\ocite{leinster}. On the other hand, there are a number of rather easy to
define examples of (weak) 3- and 4-categories, some of which we will give
later.


An alternative approach to the notion of a weak $n$-category is to use the
ideas of topology, or more precisely, homotopy theory. The key idea of this
approach is that instead of defining a single composition operation for
$k$-morphisms which is weakly associative, we allow many composition
operations, which must form a contractible space. Paths in this contractible
space correspond to $(k+1)$-isomorphisms between the corresponding
compositions, homotopies between paths correspond to $(k+2)$-isomorphisms, etc.
Different compositions give different points in this contractible space; the
fact that they are all isomorphic follows from the fact that this space is
connected, and all compatibility relations between these isomorphisms follow
from the fact that the space of compositions is contractible.

This idea can be used to define weak $n$-categories; it can also be used to
define $(\infty, n)$ categories. Informally, an $(\infty, n)$ category is a
category which has morphisms of all orders $k \ge 1$, but for $k>n$, all
$k$-morphisms are invertible.


Of course, the above is not yet a definition; even for $n=1$, developing this
idea and turning it into a precise definition is non-trivial. In fact, there are
many competing ways to define the notion of $(\infty, 1)$ category, all giving
different formalizations of the above idea. Among them:

\begin{itemize}
  \item Quasicategories, also known as weak Kan complexes
  \item Complete Segal spaces
  \item Segal categories
\end{itemize}

Fortunately, it turns out that in a certain sense all these definitions are
equivalent. There are a number of papers discussing this, including
\ocite{toen}, \ocite{bsp}, and the book \ocite{riehl-verity}; for a short
introduction, there is also Emily Riehl's article \ocite{riehl-notices} in the
\emph{Notices of the AMS}. We will not be discussing this equivalence.


\section{$(\infty, n)$ categories}
In this section we will given an overview of facts related to $(\infty, n)$
categories and functors. To keep the exposition reasonably short, it will be a
high-level description, without detail. Some detail, including a rigorous
definition of one possible model of $(\infty, n)$ categories, is given in
Appendix FIXME. Here we only give an informal description.

Informally, an $(\infty, n)$ category $\calC$ consists of class of  objects
(0-morphisms) $\calC_0$; for every pair of objects $x_0, x_1\in \calC_0$, we
have a set of $1$-morphisms $\calC_1(x_0, x_1)$; for every pair of
$1$-morphisms $f_0, f_1\in \calC(x_0,x_1)$, we have a set of $2$-morphisms
$\calC_2(f_0,f_1)$, and so on for all orders. Moreover, we have the following
facts.

\begin{enumerate}
  \item For any $k$-morphism $x$, we are given an identity  $k+1$-morphism
    $\id_x \in \calC_{k+1}(x,x)$.
  \item For any $k$-morphisms $x,y$, we have a distinguished subset of
    $\calC_{k+1}(x,y)$, called \emph{invertible} morphsisms, or isomorphisms. In
    particular, the identity morphism is invertible. We say that $x,y$ are
    isomorphic if there is an isomorphism $f\colon x\to y$. In thsi case, there
    also exists an isomorphism $g\colon y\to x$.
  \item In an $(\infty, n)$ catgeory, all $k$-morphisms with $k>n$ are
    invertible.
  \item For any topological space $X$, its fundamental groupoid $\Pi_\infty(X)$
    is an $(\infty, 0)$ category, with morphisms defined as in FIXME. Moreover,
    every $(\infty, n)$ category has the form $\Pi_\infty(X)$ for some
    topological space $X$.
  \item For any pair of objects $x,y$ in an $(\infty, n)$ category   $\calC$,
    the collection $\calC_\bullet(x,y)$ of all $1$-morphisms between $x,y$,
    together with higher morphisms between those, form an $(\infty,
    n-1)$-category: $k$-morphisms of $\calC_\bullet(x,y)$ are $k+1$-morphisms
    of $\calC$.
  \item Given two $(\infty, n)$ categories $\calC$, $\calD$, we can define
    set of functors $\Fun(\calC, \calD)$, which itself is an $(\infty, n)$
    category FIXME
  \item Symmetric monoidal $(\infty, n)$ categories.
\end{enumerate}


 we will
restrict ourselves to giving just one definition, based on Segal categories.
This construction takes some time; we will do it in the next several sections.
Our exposition follows the paper \ocite{schommer-pries-category}; a detailed
exposition of the theory of Segal categories as a model for $(\infty, n)$
categories can be found in \ocite{simpson}.




\begin{example}\label{x:poincare-segal}
  Let $X$ be a topological space. Then the fundamental groupoid
  $\Pi_{\infty}(X)$, informally described in \seref{s:fundamental-groupoid},
  can be rigorously defined as a Segal category. Namely, we let
  $X_0=\text{ set of points of $X$}$ (considered with discrete topology),
  and for $n\ge 1$, we define
  $$
    X_n= \{\text{continuous functions $\Delta^n\to X$}\}
  $$
  where $\Delta^n$ is the topological $n$-dimensional simplex \eqref{e:simplex};
  in other words, this is exactly the definition of singular complex functor
  $\Sing$, but now we consider $X_n, n\ge 1$, as a topological space rather
  than a set.

  In particular, we see that $X_1$ is the topological space of paths in $X$,
  and the Segal map $\mathbf{s}_n\colon X_n \to X_1\times_{X_0}\dots\times_{X_0}
  X_1$ sends a map $f\colon \Delta^n\to X$ to its restriction to the collection
  of edges $[i-1, i]$. Thus, one can think of a point $f\in X_n$ as a composable
  collection $\ga_i$ of paths in $X$ together with a map $\ga \colon
  \Delta^n\to X$ whose restriction to edge $[i-1, i], i=1\dots n$, is equal to
  $\ga_i$. The Segal condition immediately follows from the fact that the union of
  these edges is a retract of the $n$-simplex.

  In particular, given such a map $\ga$, we can restrict it to the edge $[0,n]$,
  which gives us a path connecting $x_0$ with $x_n$ in $X$. Any such path can be
  called a composition of $\ga_i$; thus, we see that we have not one but many
  possible compositions, but the space of all possible compositions (for given
  paths $\ga_i$) is contractible.

  Segal category defined above is usually called the \emph{Poincar\'e--Segal
  groupoid}; see \ocite{simpson}*{Sections 15.2, 15.3}. We will denote it
  $\Pi_S(X)$ to distinguish it from (informally defined) $\Pi_{\infty}$.

\end{example}


In particular:
\begin{itemize}
  \item An $(\infty, 0)$ category is the same as a topological space.
  \item An $(\infty, 1)$ category is the same as a Segal category as defined
    in \seref{s:segal-cat}.
  \item \dots
\end{itemize}

Note that by the definition of $\Seg^n$, if $\calC$ is an $(\infty, n)$
category, $n>0$, then for any two objects $x,y\in \calC$, $\Hom_\calC(x,y)$ is
an $(\infty, n-1)$ category. We can use it to define $k$-morphisms in an
$(\infty, n)$ category: for $n=0$ (i.e., when $\calC=X$ is a topological space),
$k$-morphisms are defined as in \seref{s:fundamental-groupoid}; for $n>0$,
$k$-morphisms in $\calC$ are $(k-1)$-morphisms in the $(\infty, n-1)$ category
$\Hom_{\calC}(x,y)$ for some objects $x,y\in \calC$.

We also note that the functor $\tauzero=\pi_0\colon \Top\to\sets$ gives rise to
a truncation functor $\Seg^n(\Top)\to \Seg^n(\sets)$; thus, any $(\infty, n)$
category can be truncated to a (Tamsamani) weak $n$-category. Moreover,
combining it with truncation functors $\tau_{\leq k}\colon \Tam^n\to \Tam^k$, $k
\le n$, we can define a truncation functor $\tau_{\leq k}$ from $(\infty, n)$
categories to weak $k$-categories.


%%%%%%%%%%%%%%%%%%%%%%%%%%%%%%%%%%%%%%%%%%%%%%%
\section{Symmetric monoidal infinity categories}\label{s:monoidal-infty}

FIXME

Example: let $G$ be a nice topological group; then $\Pi_\infty(G)$ is a
monoidal category, and similarly for $\Pi_{\le n}(G)$. In this case, delooping
gives us
$$
\calB(\Pi_{\le k}(G))=\Pi_{\le k+1}(BG)
$$





%%%%%%%%%%%%%%%%%%%%%%%%%%%%%%%%%%%%%%%%%%%%%%%
\section{Bordism category}\label{s:bordn-highercat}
Recall that we have defined the notion of manifolds with corners.

\begin{theorem}\label{t:bordn-highercat}
  There exists an $(\infty, n)$ category $\Bord_n$ whose objects are points,
  \dots, $n$-morphisms are $n$-dimensional manifolds with corners.
\end{theorem}
This construction was first suggested by Lurie in \ocite{lurie-cobordism}; an
improved version, fixing some shortcomings of Lurie's, is given in
\ocite{calaque}. Note, however, that both references define $(\infty, n)$
categories in terms of $n$-fold complete Segal spaces, which is slightly
different from the Segal $n$-category approach we take here. However, as
mentioned before, all models of $(\infty, n)$ categories are equivalent, so one
should be able to rewrite their construction in the language of Segal
$n$-categories (however, to the best of our knowledge, it has not been
explicitly done yet).


\clearpage